\section{The count-only question at finite and formal order}
\label{sec:countfibers}
The signed observation and the unmarked count are different experiments.
A simultaneous transverse reflection leaves both count germs unchanged.
That fact alone does not decide whether exact analytic count germs
determine an asymmetric pair modulo reflection. We give a more precise
finite-order statement, and specify what it does and does not imply at
infinite order.

\begin{proposition}[Finite count fibers]\label{prop:countfibers}
Fix the positive leading geometry and $M\ge2$. The map
\[
 (Q_3,\ldots,Q_{2M})\longmapsto(\xi_1,\ldots,\xi_{M-1})
\]
has analytic local fibers of dimension $2M-2$. The odd jets
$Q_3,Q_5,\ldots,Q_{2M-1}$ are coordinates along these fibers; the
even jets are determined successively by them and the fixed count
coefficients. Near the physical tables of Proposition~\ref{prop:realization},
these fibers occur in actual analytic periodic tables at fixed gap,
curvatures and area. Quotienting by simultaneous reflection does not
remove their positive dimension.
\end{proposition}
\begin{proof}
At stage $m$ the count equation is
\begin{equation}\label{eq:countfiberrecursion}
 \xi_{m-1}=B_{2m}Q_{2m}
             +F_{2m}(Q_3,\ldots,Q_{2m-1}),\qquad 2\le m\le M.
\end{equation}
The matrix $B_{2m}$ is \eqref{eq:evenblock}, and is invertible.
Choose all the odd jets freely in a small neighborhood. At $m=2$
solve \eqref{eq:countfiberrecursion} for $Q_4$; then solve for $Q_6$,
and continue to $Q_{2M}$. This yields an analytic graph over the
$2(M-1)$ odd variables. Conversely every solution is on this graph.
The complete jet-coordinate map in Proposition~\ref{prop:realization}
realizes each nearby solution in an analytic periodic table with the
specified leading geometry and area. A finite reflection group cannot
make a positive-dimensional local analytic fiber discrete: away from
its fixed set its quotient is locally finite-to-one, and arbitrarily
many distinct reflection orbits remain also near the fixed set.
\end{proof}

The coefficient map in this proposition is truncated. Equality of these
$M-1$ coefficients is not equality at finitely many positive windows,
whose higher-order remainders may carry additional information.
The positive-window coordinate theorem for the marked experiment below
is consequently consistent with these count fibers.

\begin{proposition}[Formal fibers and smooth realization]
\label{prop:formalcount}
At fixed leading geometry, any prescribed sequence of formal odd contact
jets determines a unique sequence of formal even jets with prescribed
formal count coefficients, by \eqref{eq:countfiberrecursion}. Near a
fixed strictly dispersing periodic reference geometry, such compatible
formal contact jets can be realized by smooth physical tables with the
same contact locations, gap, curvatures and free area. In particular,
there are pairs not related by common reflection whose normalized count
laws have the same Taylor series at onset.
\end{proposition}
\begin{proof}
The formal assertion follows by induction, since each equation uses only
finitely many earlier jets and the same nonsingular even block. There
is no infinite equation to solve at any induction step.

For the physical assertion one may start with the unit-disk construction
of Proposition~\ref{prop:realization}. The two registered contacts are
distinct points of one boundary modulo the lattice. A prescribed graph
jet there determines, degree by degree, a support-function jet: at each
degree the diagonal support-to-graph derivative is the nonzero value
in \eqref{eq:supportjac}. The leading support value, first derivative
and curvature radius are kept fixed. We recall explicitly the needed
smooth jet extension. In disjoint support-angle neighborhoods of the
two contacts, multiply the desired degree-$r$ Taylor monomial by a
smooth cutoff equal to one near its origin, with radius $\varepsilon_r$.
Choose the radii decreasing so rapidly that the $r$th summand has
$C^k$ norm at most $2^{-r}\epsilon$ for every $k\le r-1$.
This is possible because the degree-$r$ monomial has such norms of
order $\varepsilon_r^{r-k}$, with fixed coefficients at that stage.
The resulting series converges in every fixed $C^k$ norm after omitting
finitely many terms. Since each cutoff is identically one near its
origin, it has exactly the prescribed jets. Beginning at degree three
makes the full perturbation arbitrarily small in $C^2$, regardless of
the later formal coefficients.

Choose, in addition, a nonnegative nonzero smooth support perturbation
$f$ away from both contacts. It changes no contact jet. The free-area
derivative in this direction is
$-\int(h+h'')f\dd\theta<0$. The implicit-function theorem therefore
restores the reference area by a small multiple of $f$. All support
perturbations can be kept sufficiently small in $C^2$ that
$h+h''>0$, embedding, obstacle separation and channel clearance persist.
These are actual smooth periodic billiards. Finite-jet dependence of
the coefficient functionals identifies their count Taylor coefficients
with the prescribed formal ones.

To obtain the last assertion, prescribe the same count series as the
reference table but choose a nonzero cubic jet, or two odd sequences
not related by common negation, before solving the even recursion.
The resulting smooth jets, and hence the boundary germs, are not a
common reflection pair. Their count Taylor series are identical.
\end{proof}

\begin{remark}[Exact analytic count germs]\label{rem:analyticcount}
No convergence assertion for the formal inverse in
Proposition~\ref{prop:formalcount} is used or established. Its smooth
realizations may have count laws differing by a nonzero flat germ.
Thus the proposition neither proves nor disproves determination by
\emph{exact analytic} count germs modulo common reflection. It shows why
finite-order counting and the reflection obstruction alone cannot settle
that question. Establishing convergence of suitable formal fibers, or a
rigidity mechanism that excludes them in the analytic class, would be
needed to decide it. We do not assign either outcome on the basis of a
formal recursion. The intrinsic marked theorem avoids this unresolved
step by determining each odd jet directly with a nonsingular block.
\end{remark}
